\clearpage\thispagestyle{empty}
\pdfbookmark[0]{Sources for the 10 March 1807 Letter}{zach18070310-sources}
\begin{center}\textbf{Sources for the 10 March 1807 Letter}\end{center}
\noindent This AI translation follows Carl Friedrich Gauss’s letter of 10 March 1807 to Franz Xaver von Zach in \href{https://gdz.sub.uni-goettingen.de/id/PPN236007467}{Gauss’s \textit{Werke}, vol.~VI} (Göttingen, 1874), pp.~281--282. It appeared as article XXIV in the April 1807 \textit{Monatliche Correspondenz}, vol.~XV, \href{https://zs.thulb.uni-jena.de/rsc/viewer/jportal_derivate_00237836/Monatliche_Correspondenz_130168688_15_1807_0380.tif}{pp.~377}--\href{https://zs.thulb.uni-jena.de/rsc/viewer/jportal_derivate_00237836/Monatliche_Correspondenz_130168688_15_1807_0381.tif}{378}. The journal scans are from the Thüringer Universitäts- und Landesbibliothek Jena.
\par\medskip
On \href{https://archive.org/details/bub_gb_sKAAAAAAMAAJ/page/n164/mode/1up}{p.~159 of the February issue}, in the observations sent by Oriani, the declination Carlini gives for 4 January 1806 is $29^\circ\,20'\,29''3$. Gauss gives $29^\circ\,21'\,29''3$ in his letter and discusses a separate error in latitude. The February issue is from the New York Public Library copy digitized by Google and available through the Internet Archive.
\par\medskip
The abbreviation \emph{Hist. cél.} refers to Jérôme de Lalande’s \href{https://archive.org/download/histoireclestef00lalagoog/page/n12.jpg}{\textit{Histoire céleste française}} (Paris, An~IX [1801]). Translation and typesetting produced by AI.
\clearpage
